\section{Fatiamento de Redes}

Uma fatia de rede é uma rede virtual implementada a partir dos recursos de uma infraestrutura física. Por meio do fatiamento de redes, uma mesma infraestrutura pode ser utilizada para criar diferentes fatias, permitindo que cada uma seja configurada de acordo com requisitos específicos de serviços ou aplicações \cite{zhang2017slicingbased5g}. Dessa forma, diferentes fatias podem compartilhar os recursos da infraestrutura física, ao mesmo tempo em que são estabelecidas para atender diferentes requisitos de comunicação.

O conceito de fatiamento de redes está relacionado a abordagens de virtualização que permitem o compartilhamento de uma infraestrutura física entre diferentes redes virtuais. Entretanto, o fatiamento ganhou maior destaque no contexto das redes móveis 5G, diante da necessidade de uma mesma infraestrutura suportar diferentes serviços e aplicações com requisitos distintos de desempenho e comunicação \cite{foukas2017network}.

O fatiamento de redes tem como objetivo ir além do compartilhamento de recursos ou da redução de custos, buscando fornecer redes e serviços customizados de acordo com diferentes requisitos. Para isso, são necessários modelos e mecanismos mais flexíveis para o provisionamento e a implantação dos serviços de rede \cite{boubendir2019network}. Essa flexibilidade permite que diferentes fatias sejam configuradas para atender às necessidades específicas dos serviços e aplicações que utilizam a infraestrutura.

A infraestrutura dos ISPs é compartilhada entre usuários e serviços que podem apresentar diferentes necessidades de recursos ao longo do tempo. Nesse contexto, o fatiamento de redes possibilita dividir logicamente os recursos da infraestrutura em diferentes fatias, permitindo que sejam alocados de acordo com os requisitos de cada serviço. Dessa forma, a capacidade disponível pode ser distribuída entre as fatias e adaptada às diferentes condições de utilização da rede \cite{gomes2020strategies}.

Entretanto, a definição de como os recursos devem ser distribuídos entre as diferentes fatias constitui um dos principais desafios do fatiamento de redes, especialmente em cenários nos quais a demanda por recursos varia ao longo do tempo. A existência de diferentes requisitos, como largura de banda, latência e disponibilidade, torna necessário o desenvolvimento de mecanismos capazes de realizar essa alocação de forma eficiente. Nesse cenário, os algoritmos de fatiamento desempenham um papel importante, uma vez que auxiliam na determinação de como os recursos disponíveis serão alocados entre as diferentes fatias \cite{vassilaras2017algorithmic}.

Além disso, a alocação dos recursos pode precisar ser adaptada às novas condições da infraestrutura, de acordo com as variações da demanda ao longo do tempo. Nesse cenário, as necessidades de uma fatia podem aumentar ou diminuir, exigindo alterações na quantidade de recursos alocados ou até mesmo em sua configuração. Esse comportamento está relacionado ao fatiamento dinâmico, no qual as fatias podem ser reconfiguradas para acompanhar as mudanças nas condições da rede \cite{vassilaras2017algorithmic}.

Além da flexibilidade na alocação e adaptação dos recursos, o fatiamento de redes apresenta características que favorecem o gerenciamento e a utilização da infraestrutura. \citeonline{vassilaras2017algorithmic} destacam, entre seus benefícios, o isolamento entre fatias, a simplificação das cadeias de serviços, a flexibilidade no posicionamento de funções virtuais de rede e a transparência no gerenciamento das fatias:
\begin{itemize}
    \item \textbf{Isolamento de fatias:} o isolamento permite que cada fatia opere de maneira independente das demais, mesmo quando compartilham recursos de uma mesma infraestrutura física. Dessa forma, falhas de rede ou problemas de segurança ocorridos em uma fatia podem ser isolados, reduzindo seus impactos sobre as demais fatias.
    \item \textbf{Cadeias de serviços simplificadas:} diferentemente de arquiteturas nas quais um conjunto amplo de funções de rede pode estar presente mesmo quando nem todas são necessárias, o fatiamento permite que cada fatia seja composta pelas funções adequadas aos seus requisitos. Dessa forma, as cadeias de serviços podem ser simplificadas, utilizando apenas as funções necessárias para os serviços oferecidos pela fatia.
    \item \textbf{Posicionamento flexível de VNFs:} a possibilidade de posicionar as Funções Virtuais de Rede (VNFs) em diferentes locais da infraestrutura adiciona maior flexibilidade à composição das fatias. Um posicionamento adequado dessas funções pode contribuir para melhorar o desempenho da rede e reduzir os recursos necessários para o fornecimento dos serviços.
    \item \textbf{Gerenciamento transparente de fatias:} os recursos físicos utilizados por uma fatia podem pertencer a diferentes domínios, o que pode aumentar a complexidade de seu gerenciamento. Entretanto, o fatiamento de redes fornece uma abstração desses recursos, tornando seu gerenciamento mais transparente e facilitando o tratamento da fatia independentemente da localização dos recursos físicos utilizados.
\end{itemize}

Diante dessas características, o fatiamento de redes apresenta-se como uma abordagem capaz de contribuir para o gerenciamento de recursos em ambientes de demanda elástica. A possibilidade de criar fatias com diferentes configurações e adaptar a utilização dos recursos de acordo com seus requisitos permite que a infraestrutura seja utilizada de forma mais flexível. Nesse contexto, diferentes técnicas e estratégias de fatiamento podem ser empregadas para determinar como os recursos disponíveis serão distribuídos e adaptados às necessidades de cada fatia.